\documentclass[10pt,a4paper]{amsart}
\usepackage[margin=19mm]{geometry}
\usepackage{amsmath,amssymb,mathtools}
\usepackage[scr=rsfs,cal=euler]{mathalfa}
\usepackage{xfrac}
\usepackage[unicode,hidelinks,bookmarks=false]{hyperref}
\newcommand{\ie}{i.e.\ }
\newcommand{\cf}{cf.\ }
\newcommand{\resp}{respectively\ }
\newcommand{\Subsection}{\subsection}
\providecommand{\nobreakdash}{\nobreak\hbox{-}}
\setlength{\emergencystretch}{2em}
\multlinegap0pt
\usepackage{amssymb}
\usepackage{enumerate}
\usepackage{xcolor}
\newtheorem{lemma}{Lemma}
\newtheorem{theorem}{Theorem}
\title{Sections and cones}
\author{Tatiana Shulman}
\date{Source: arXiv:2507.22783v3 (CC BY 4.0)}
\begin{document}
\maketitle

\noindent\emph{Source and licence.} Sections and cones. Version arXiv:2507.22783v3 (CC BY 4.0), distributed by arXiv under the Creative Commons Attribution 4.0 International licence, http://creativecommons.org/licenses/by/4.0/, as recorded in the arXiv licence metadata for this exact version and shown on the abstract page https://arxiv.org/abs/2507.22783v3. Full licence notice: \texttt{licenses/arxiv-2507.22783-CC-BY-4.0.txt}.\par\smallskip

\noindent\textit{Editorial scope.} Counted unit: the article's theorem ccp (source0.tex lines 779--781). The article's complete original proof of it is delivered with it (source0.tex lines 782--810, one \texttt{proof} environment of 29 lines). No mathematics is written, repaired or re-derived: the statement and the proof are the article's own lines, in the article's own notation, and no retained line is substituted or re-flowed.  Retained uncounted as the source-local context the proof needs: lines 480--482; lines 551--552; lines 555--556; lines 718--719; lines 764--769.  Nothing was omitted from any retained span, so the package carries no omission record.  No retained line contains a citation, so the wrapper carries no bibliography and no bibliography entry.  No retained line contains a figure environment, an image-inclusion command or drawing code, so no image is delivered and the package declares no asset dependency.  Editorial formatting only, in addition: an AMS \texttt{amsart} preamble that restates the article's own declarations and macros used by the retained lines, namely the environments \texttt{lemma}, \texttt{theorem} on separate counters, together with the standard packages those lines need.  The retained environments print in this document as Lemma 1, Theorem 1 in order of appearance, whereas the article prints the same items with the numbering of its own full document.  The complete native source of the article as distributed in the arXiv e-print for this exact version is bundled unchanged as \texttt{sources/182441/source0.tex}.
\begin{lemma}\label{AsHom1}  Let $p (x_1, \ldots, x_N) $ be a homogeneous NC $\ast$-polynomial (more generally, $p$ can be of more variables and homogeneous only in $x_1, \ldots, x_N$). Let $I \lhd B$, $\{i_{\lambda}\}$ a quasicientral approximate unit of $I$ relative to $B$, and $\pi: B \to B/I$  the canonical surjection.
 Suppose $\pi(p(b_1, \ldots, b_N)) = 0$. Then $$\lim \|p(b_1(1-i_{\lambda}), \ldots, b_N(1-i_{\lambda}))\| =0.$$
\end{lemma}

\begin{lemma}\label{continuousqau} Let $u_n, n\in \mathbb N,$ be a quasicentral approximate unit for $I\triangleleft B$. Let $u_t = (n+1-t)u_n + (t-n)u_n$, for $n\le t\le n+1$, $n\in \mathbb N$. Then the continuous path $u_t$, $t\in [1, \infty)$, is also a quasicentral approximate unit for $I\triangleleft B$.
\end{lemma}

\begin{theorem}\label{cones} Let $A$ be a  separable C*-algebra. Then any $\ast$-homomorphism from $CA$ to any quotient $B/I$ lifts to a contractive positive asymptotic homomorphism.
\end{theorem}

\begin{theorem}\label{LiftingDuality}   Let $A$  be a separable C*-algebra, $X$ a locally compact space, $\xi\in C_0(X)$ a strictly positive element, and let $\delta: A \to A\otimes C_0(X)$ be defined by $\delta(a)= a\otimes \xi$, $a\in A$. Let    $\rho: A\otimes C_0(X)\to B/I$  be a cp map.  Then $\rho$  lifts to a cp map if and only $\rho\circ\delta$ lifts to a cp map.
\end{theorem}

\begin{lemma}\label{cpLiftFromUnitization} \begin{itemize}
 \item[(1)] Let $0\to J\to E\to A\to 0$ be a  short exact sequence, where $E$ is separable  and  $A$ has the  Lifting Property (LP). Let  $\psi: E\to B/I$ be a cp map. If $\psi|_J$ has a cp lift, then $\psi$ has a cp lift.

\item[(2)] If $\psi: C(A^+)\to B/I$ is a cp map and $\psi|_{CA}$ has  a cp lift, then $\psi$ has  a cp lift.
\end{itemize}
\end{lemma}

\begin{theorem}\label{ccp}   Let $A$ be separable and let $f: CA \to B/I$ be a $\ast$-homomorphism that has a cpc lift $\phi: CA \to B$.
 Then $f$ lifts to a cpc  asymptotic homomorphism.
  \end{theorem}

  \begin{proof} If $A$ is non-unital, we extend $f$ to $\tilde f: C(A^+)\to M(B/I)$ in the same way as in part 2) of the proof of Theorem \ref{cones}.
  By Lemma \ref{cpLiftFromUnitization}, $\tilde f$ lifts to a cpc map $ \psi: C(A^+)\to M(B)$. If $A$ is unital, we let $\psi = \phi$.



  Let $\{i_{\lambda}\}_{\lambda\in \Lambda}$ be  a quasicentral approximate unit of $I$ relative to $B$. Here $\Lambda=\mathbb N$ is we are interested in discrete asymptotic homomorphisms and $\Lambda = [1, \infty)$ for the continuous version in which case we take a quasicentral approximate unit forming a continuous path as in Lemma \ref{continuousqau}.
Let $$\psi_{\lambda} = (1-i_{\lambda})^{1/2} \psi (1-i_{\lambda})^{1/2}: C(A^+)\to M(B).$$ For each $\lambda\in \Lambda$ and $x\in CA$ we have $$q(\psi_{\lambda}(x)) = f(x).$$
 Let  $ a_1, a_2 \ldots$ be a dense subset of $A$. Let $x_i = \delta(a_i)$, $h = \delta(1_{A^+})$.
We write $C(A^+)$ as the universal C*-algebra with generators $h$ and $x_i$, $i\in \mathbb N$, and homogeneous relations.   By Lemma \ref{AsHom1}  $\psi_{\lambda}(h)$ and $\psi_{\lambda}(x_i)$'s approximately satisfy the relations.  Let $$\pi: \prod M(B) \to \prod M(B) /\oplus I$$ ($\pi: C_b([1, \infty), M(B))\to C_b([1, \infty), M(B))/C_0([1, \infty), I)$ for the continuous version) be the canonical surjection. Then $\pi\left((\psi_{\lambda}(h))_{\lambda\in \Lambda}\right)$ and  $\pi\left((\psi_{\lambda}(x_i))_{\lambda\in \Lambda}\right)$, $i\in \mathbb N$, satisfy all the relations of $C(A^+)$ and therefore define a $\ast$-homomorphism
$$\Psi: C(A^+) \to \prod M(B)/\oplus I$$ ($\Psi: C(A^+) \to C_b([1, \infty), M(B))/C_0([1, \infty), I)$ for the continuous version). Let
$$\Phi = \Psi|_{CA}.$$ We notice that it lands in $\prod B/\oplus I$ (in $C_b([1, \infty), B)/C_0([1, \infty), I)$ for the continuous version). We have
$$\Phi(\delta(a_i)) = \Phi(x_i) = \pi\left((\psi_{\lambda}(x_i))_{\lambda\in \Lambda}\right) =  \pi\left((\psi_{\lambda}(\delta(a_i)))_{\lambda\in \Lambda}\right). $$ Since $\{a_i\}$ is dense in $A$, we conclude that
$$\Phi\circ \delta =\pi\circ \left(\psi_{\lambda}|_{CA}\circ \delta\right)_{\lambda\in \Lambda}.$$
Thus $\Phi\circ \delta$ has a cp lift $\left(\psi_{\lambda}|_{CA}\circ \delta\right)_{\lambda\in \Lambda}$. (We note that one cannot say that $\left(\psi_{\lambda}|_{CA}\right)_{\lambda\in \Lambda}$ is a cp lift of $\Phi$. It only lifts $\Phi$ on the linear span of the generators). Then, by Theorem \ref{LiftingDuality}, $\Phi$ has a cp lift
 $$\left(f_{\lambda}\right)_{\lambda\in\Lambda}: CA \to \prod B$$ ($\left(f_{\lambda}\right)_{\lambda\in\Lambda}: CA \to C_b([1, \infty, B)$ for the continuous version. In this case we also obtain that the map $\lambda\mapsto f_{\lambda}$ is continuous).



  Then for any $x, y\in CA$ we have
 $$\left(f_{\lambda}(xy) - f_{\lambda}(x)f_{\lambda}(y)\right)_{\lambda\in \Lambda}\in \oplus I  \; $$
  ($\in (C_0([1, \infty), I), \;\text{respectively}$) and similar for linear combinations and adjoint elements. Therefore
 $f_{\lambda}$, $\lambda\in \Lambda$, is an asymptotic homomorphism such that
  $q\circ f_{\lambda}$ is a $\ast$-homomorphism for each $\lambda\in \Lambda$.

  By the construction of $\Phi$ and $f_{\lambda}$, for each generator $x_i$ of $CA$ we have
  $$\pi\left((f_{\lambda}(x_i))_{\lambda\in\Lambda}\right) = \Phi(x_i) = \pi\left((\psi_{\lambda}(x_i))_{\lambda\in\Lambda}\right).$$
  Hence $$q\circ f_{\lambda}(x_i) = q\circ \phi_{\lambda}(x_i) = f(x_i),$$ for each $\lambda\in \Lambda$. Therefore
  $q\circ f_{\lambda} = f, $ for each $\lambda \in \Lambda$.
  \end{proof}

\end{document}
